\documentclass[10pt,compress,t,notes=noshow, xcolor=table]{beamer}

\input{../../style/preamble}
\input{../../latex-math/basic-math}
\input{../../latex-math/basic-ml}
\input{../../latex-math/ml-interpretable}
% Defines macros and environments

\title{Interpretable Machine Learning}
% \author{LMU}
%\institute{\href{https://compstat-lmu.github.io/lecture_iml/}{compstat-lmu.github.io/lecture\_iml}}

\newcommand{\gh}{\hat{g}}

\begin{document}

% Set style/preamble.Rnw as parent.

% Load all R packages and set up knitr

% This file loads R packages, configures knitr options and sets preamble.Rnw as 
% parent file
% IF YOU MODIFY THIS, PLZ ALSO MODIFY setup.Rmd ACCORDINGLY...

% Defines macros and environments

\titlemeta{
Local Explanations: Lime
}{
Local Interpretable Model-agnostic Explanations (LIME)
}{
figure/lime5
}{
\item Understand motivation for LIME
\item Develop a mathematical intuition
}


% Prerequisite: le-intro

% ------------------------------------------------------------------------------
\begin{framei}[align=top]{LIME}
\item \textbf{Locality assumption:}\\
$\fh$ behaves similarly simple in a small neighborhood of $\xv$\\
$\leadsto$ Approximate $\fh$ near $\xv$ using an interpretable surrogate model $\gh$
\item \textbf{Interpretation strategy:}\\
Use $\gh$'s simple internal structure to explain $\fh(\xv)$ locally\\
$\leadsto$ \textbf{Common surrogates:} sparse linear models, shallow decision trees
\item \textbf{Applicability:} Model‑agnostic; supports tabular, image, and text data %via modality‑specific perturbations
\item \textbf{In practice:}
Generate samples near $\xv$, predict with $\fh$, and fit $\gh$ to these samples using $\fh$'s outputs as targets, weighting samples by their proximity to $\xv$
\end{framei}


\begin{framev}[align=top]{LIME: Characteristics}
\textbf{Definition:}\\
LIME provides a local explanation for a black-box model $\fh$ in form of a surrogate model $\gh \in \Gspace$, where $\Gspace$ is a class of interpretable models\\
\spacer[0.25]
Surrogate model $\gh$ should satisfy two characteristics:
\begin{enumerate}
\item \textbf{Interpretable}: Provide human-understandable insights into the relationship between input features and prediction, e.g. via coefficients or model structure
\item \textbf{Local fidelity / faithfulness}: $\gh$ closely approximates $\fh$ in the vicinity of the input $\xv$ being explained
\end{enumerate}
\textbf{Goal:} Find $\gh$ with \textbf{minimal complexity and maximal local fidelity}
\end{framev}


\begin{framev}[align=top]{Model Complexity}
We can measure complexity of $\gh \in \Gspace$ using a complexity measure $J: \mathcal{G} \to \R_0$ \lz % $J(\gh)$ \lz
\begin{itemizeM}
\item \textbf{Example: (Sparse) Linear Models}\\
Let $\Gspace = \{g: \Xspace \to \R ~|~ g(\xv) = s(\thetav^\top \xv)\}$ be the class of linear models
\begin{itemizeS}
\item $s(\cdot)$ is the identity (linear model) or logistic sigmoid function (logistic regression)
\item[$\leadsto$] $J(g) = \sum_{j = 1}^p \mathds{1}_{\{\theta_j \neq 0\}}$: count the number of non-zero coefficients via the $L_0$ norm of $\thetav$
\end{itemizeS}
\pause
\item \textbf{Example: Decision Trees}\\
Let $\Gspace = \{g: \Xspace \to \R ~|~ g(\xv) = \sum_{m=1}^M c_m \mathds{1}_{\{\xv \in \mathcal{R}_m\}}\}$ be the class of trees
\begin{itemizeS}
\item $\mathcal{R}_m$ are disjoint axis-parallel regions (leaves); $c_m \in \R$ are constant predictions
%\\ 	     i.e., the class of additive models (e.g., constant $c_m$)  over the leaf-rectangles $\mathcal{R}_m$
\item[$\leadsto$] $J(g) = M$: count the number of terminal or leaf nodes
\end{itemizeS}
\end{itemizeM}
\end{framev}
 
\begin{framev}[align=top]{Local Fidelity of Surrogate Models}
\begin{itemizeM}
\item Surrogate $\gh$ is \textbf{locally faithful} to a black-box model $\fh$ around an input $\xv$ if\\
$$\gh(\zv) \approx \fh(\zv) \quad \text{for synthetic samples } \zv \in \Zspace \subseteq \R^p \text{ generated around } \xv$$
\pause
\item \textbf{Optimization principle:}
The closer $\zv$ is to $\xv$, the more $\gh(\zv)$ should match $\fh(\zv)$
\pause
\item To operationalize this optimization, we need:
\begin{enumerate}
\item \textbf{A proximity (similarity) measure} $\neigh(\zv)$ between $\zv$ and $\xv$, e.g.:
$$
\neigh(\zv) = \exp(-d(\xv, \zv)^2 / \sigma^2)
$$
(exponential kernel), where
\begin{itemizeS}
\item $d$ is a distance metric, e.g. Euclidean or Gower for mixed types
\item $\sigma$ is the kernel width that controls locality
\end{itemizeS}
%(e.g., Euclidean or Gower),
\pause
\item \textbf{A loss function} $L(\fh(\zv), \gh(\zv))$, e.g. the $L_2$ loss or squared error:
$$
L(\fh(\zv), \gh(\zv)) = (\gh(\zv) - \fh(\zv))^2
$$
\end{enumerate}
\pause
\item The overall \textbf{local fidelity objective} is measured by a weighted loss:
$$
L(\fh, \gh, \neigh) = \sum_{\zv \in \Zspace} \neigh(\zv) \cdot L(\fh(\zv), \gh(\zv))
$$
%\item Note that identifying \textbf{locally} faithful explanations that are interpretable is less of a challenge than identifying \textbf{globally} faithful explanations. Yet, global fidelity implies local fidelity but not vice versa.
\end{itemizeM}
\end{framev}


\begin{framei}[align=top]{LIME Optimization Task}
\item Optimization problem of LIME:
$$
\argmin_{\gh \in \Gspace} L(\fh, \gh, \neigh) + J(\gh)
$$
\item \textbf{In practice} LIME uses a two-stage approach:
\begin{enumerate}
\item User sets complexity $J(\gh)$ beforehand, e.g. LASSO with $k$ features
\item Optimize $L(\fh, \gh, \neigh)$ for fixed complexity
\end{enumerate}
\item \textbf{Goal:} Build a \textbf{model-agnostic} explainer
\begin{itemizeS}
\item[$\leadsto$] Optimize $L(\fh, \gh, \neigh)$ without making assumptions on the form of $\fh$
\item[$\leadsto$] Surrogate $\gh$ approximates $\fh$ locally through sampling and fitting
\end{itemizeS}
\end{framei}


\begin{framev}[align=top]{LIME Algorithm: Outline \furtherreading{RIBEIRO2016}}
\textbf{Input}:
\begin{itemizeM}
\item Pre-trained black-box model $\fh$
\item Observation $\xv$ whose prediction $\fh(\xv)$ we want to explain
\item Interpretable model class $\Gspace$ for the local surrogate
\end{itemizeM}
\pause
\textbf{Algorithm}:
\begin{enumerate}
\item Independently sample new points $\zv \in \Zspace$
\item Retrieve predictions $\fh(\zv)$ for the obtained points $\zv$
\item Weight $\zv \in \Zspace$ by their proximity $\neigh(\zv)$ to quantify closeness to $\xv$
\item Train interpretable surrogate model $\gh$ on data $\zv \in \Zspace$ using weights $\neigh(\zv)$\\
$\leadsto$ Predictions $\fh(\zv)$ are used as the target of this model
\item Return $\gh$ as the local explanation for $\fh(\xv)$
\end{enumerate}
\end{framev}


\begin{framev}[align=top]{LIME Algorithm: Example}
\textbf{Illustration} of LIME based on a classification task:
\begin{itemizeM}
\item Light/dark gray background: prediction surface of a classifier
\item Yellow point: $\xv$ to be explained
\item $\Gspace$: class of logistic regression models
\end{itemizeM}
\imageC[0.6]{figure/lime2}
\end{framev}


\begin{framev}[align=top]{LIME Algo.: Example (Step 1+2: Sampling)}
Strategies for sampling:
\begin{itemizeM}
\item Uniformly sample new points from the feasible feature range
\item Use the training data set with or without perturbations
\item Draw samples from the estimated univariate distribution of each feature
\item Create an equidistant grid over the supported feature range
\end{itemizeM}
\splitVTT{%
\spacer[0]
\begin{figure}
\centering
\image[0.95]{figure/lime3}
\caption{Uniformly sampled}
\end{figure}
}{%
\spacer[0]
\begin{figure}
\centering
\image[0.95]{figure/lime3a}
\caption{Equidistant grid}
\end{figure}
}
\end{framev}


\begin{framev}[align=top]{LIME Algo.: Example (Step 3: Proximity)}
In this example, we use the exponential kernel defined on the Euclidean distance $d$:
$$
\neigh(\zv) = \exp(-d(\xv, \zv)^2 / \sigma^2)
$$
\imageC[0.7]{figure/lime4}

% MARIUS: Not relevant for the example; if we want to introduce it, we should do it somewhere else
% 	An alternative is the Gower proximity: 
% 	$\neigh(\zv) = 1 - \Gower(\zv, \xv) =  1 - \frac{1}{p}\sum_{j = 1}^{p} \delta_G(z_j, x_j) $ 
% 	$\textnormal{ with } \delta_G(z_j, x_j) = 
% 	\begin{cases}
% 	\frac{1}{\widehat{R}_j}|z_j- x_j| & \text{if $x_j$ and $z_j$ are numerical} \\
% 	\mathbb{I}_{z_j \neq x_j} & \text{if $x_j$ and $z_j$ are categorical}
% 	\end{cases}.$

\end{framev}


\begin{framev}[align=top]{LIME Algo.: Example (Step 4: Surrogate)}
In this example, we fit a \textbf{logistic regression} model\\
$\leadsto$ $L(\fh(\zv), \gh(\zv))$ is the Bernoulli loss
\imageC[0.7]{figure/lime5}
\end{framev}

\endlecture
\end{document}
